\chapter{Moments and Transforms}
\label{ch:transforms}

A distribution can be encoded in a function whose algebra matches the algebra of
sums of independent variables. Three such encodings are used throughout probability:
the moment generating function, the probability generating function (for
integer-valued variables), and the characteristic function (which always exists).
They turn convolutions into products, make moments computable by
differentiation, and are the main tool in the proof of the central limit theorem.

\section{Moments}

\begin{definition}[Moments]
For $k\ge1$ the $k$th \emph{moment} of $X$ is $\Exp[X^k]$ and the $k$th
\emph{central moment} is $\Exp[(X-\mu)^k]$, provided $\Exp\abs X^k<\infty$. The
\emph{skewness} is $\Exp[(X-\mu)^3]/\sigma^3$ and the \emph{kurtosis} is
$\Exp[(X-\mu)^4]/\sigma^4$.
\end{definition}

\begin{proposition}[Lower moments from higher ones]\label{prop:lyapunov}
If $\Exp\abs X^s<\infty$ and $0<r<s$, then $\Exp\abs X^r<\infty$.
\end{proposition}
\begin{proof}
For every real $x$, $\abs x^r\le1+\abs x^s$ (if $\abs x\le1$ then $\abs x^r\le1$;
otherwise $\abs x^r\le\abs x^s$). Take expectations.
\end{proof}

The converse is false: the distribution with density $\propto x^{-4}$ on $x>1$ has
a finite mean and variance but infinite third moment.

\section{The moment generating function}

\begin{definition}[MGF]
The \emph{moment generating function} of $X$ is $M_X(t)=\Exp[e^{tX}]$ for those
real $t$ for which the expectation is finite. We assume it is finite on an open
interval around $t=0$ whenever we use it.
\end{definition}

\begin{theorem}[Moments from the MGF]\label{thm:mgfmoments}
If $M_X(t)<\infty$ for $\abs t<\delta$ for some $\delta>0$, then all moments exist and
$M_X$ is infinitely differentiable on $(-\delta,\delta)$, with
\[
\Exp[X^k]=M_X^{(k)}(0),\qquad
M_X(t)=\sum_{k=0}^\infty\frac{\Exp[X^k]}{k!}t^k\quad(\abs t<\delta).
\]
\end{theorem}
\begin{proof}
Fix $0<t<\delta$. Finiteness of $\Exp[e^{tX}]$ and $\Exp[e^{-tX}]$ implies
$\Exp[e^{t\abs X}]\le\Exp[e^{tX}]+\Exp[e^{-tX}]<\infty$, because
$e^{t\abs x}\le e^{tx}+e^{-tx}$. Since $e^{t\abs x}\ge(t\abs x)^k/k!$ for every $k$,
we get $\Exp\abs X^k\le k!\,\Exp[e^{t\abs X}]/t^k<\infty$.
Expanding $e^{tX}=\sum_k(tX)^k/k!$ and interchanging expectation with the sum (justified
by dominated convergence with the bound $\sum_k\abs{tX}^k/k!=e^{\abs{tX}}$, which has
finite expectation; see Chapter~\ref{ch:measure}) gives the power series, whose
coefficients are $M_X^{(k)}(0)/k!$.
\end{proof}

\begin{theorem}[Properties]\label{thm:mgfprops}
\begin{enumerate}[label=(\roman*)]
\item $M_{aX+b}(t)=e^{bt}M_X(at)$.
\item If $X,Y$ are independent, $M_{X+Y}(t)=M_X(t)M_Y(t)$.
\item (Uniqueness, cited) If $M_X=M_Y<\infty$ on an interval around $0$, then $X$ and
$Y$ have the same distribution.
\end{enumerate}
\end{theorem}
\begin{proof}
(i) $\Exp[e^{t(aX+b)}]=e^{bt}\Exp[e^{(at)X}]$. (ii) $e^{tX}$ and $e^{tY}$ are
independent (functions of independent variables), so by
Theorem~\ref{thm:prodexp} the expectation of the product is the product of
expectations. (iii) is a deep theorem; it follows from the uniqueness of the
Laplace transform (or from Fourier inversion for the characteristic function in
Section~\ref{sec:cf}), and we use it without proof.
\end{proof}

\subsection*{Computations}

\begin{center}
\begin{tabular}{@{}lll@{}}
\toprule
Distribution & MGF & Domain\\
\midrule
Bernoulli$(p)$ & $q+pe^t$ & $\R$\\
Binomial$(n,p)$ & $(q+pe^t)^n$ & $\R$\\
Poisson$(\lambda)$ & $\exp\bigl(\lambda(e^t-1)\bigr)$ & $\R$\\
Geometric$(p)$ & $\dfrac{pe^t}{1-qe^t}$ & $t<-\ln q$\\[0.5em]
Exponential$(\lambda)$ & $\dfrac\lambda{\lambda-t}$ & $t<\lambda$\\[0.5em]
Gamma$(\alpha,\lambda)$ & $\Bigl(\dfrac\lambda{\lambda-t}\Bigr)^\alpha$ & $t<\lambda$\\[0.5em]
$N(\mu,\sigma^2)$ & $\exp\bigl(\mu t+\tfrac12\sigma^2t^2\bigr)$ & $\R$\\
\bottomrule
\end{tabular}
\end{center}

Derivations. Poisson: $\sum_ke^{tk}e^{-\lambda}\lambda^k/k!=e^{-\lambda}\exp(\lambda e^t)$.
Exponential: $\int_0^\infty e^{tx}\lambda e^{-\lambda x}dx=\lambda/(\lambda-t)$ for $t<\lambda$.
Gamma: the substitution $u=(\lambda-t)x$ in the gamma integral. For the standard
normal, complete the square, $tz-z^2/2=-(z-t)^2/2+t^2/2$:
\[
M_Z(t)=\frac1{\sqrt{2\pi}}\int e^{tz-z^2/2}dz=e^{t^2/2}\cdot\frac1{\sqrt{2\pi}}\int e^{-(z-t)^2/2}dz=e^{t^2/2},
\]
and then $M_{\mu+\sigma Z}(t)=e^{\mu t}M_Z(\sigma t)$ by (i).

\begin{example}[Sums revisited]
MGFs give instant proofs of the closure theorems. If $X\sim\mathrm{Poisson}(\lambda)$,
$Y\sim\mathrm{Poisson}(\mu)$ are independent then
$M_{X+Y}=\exp((\lambda+\mu)(e^t-1))$, so $X+Y\sim\mathrm{Poisson}(\lambda+\mu)$ by uniqueness.
Likewise, normals add by adding exponents; and a sum of $n$ independent
Exp$(\lambda)$ has MGF $(\lambda/(\lambda-t))^n$, the Gamma$(n,\lambda)$ MGF.
\end{example}

\begin{example}[Moments of the normal]
Since $M_Z(t)=e^{t^2/2}=\sum_m\frac{t^{2m}}{2^mm!}$, matching with
$\sum_k\Exp[Z^k]t^k/k!$ gives $\Exp[Z^{2m}]=\frac{(2m)!}{2^mm!}=(2m-1)!!$ and all
odd moments zero. So $\Exp[Z^4]=3$ (kurtosis $3$, the benchmark for ``normal
tails'').
\end{example}

\begin{nuance}{the MGF may not exist}
Heavy-tailed distributions have no MGF: the log-normal and Cauchy distributions
have $\Exp[e^{tX}]=\infty$ for every $t>0$. Moreover, having all moments finite is
\emph{not} enough for the moments to determine the distribution: the log-normal
distribution has finite moments of every order but shares them with other
distributions. Existence of the MGF in a neighbourhood of $0$ is a sufficient
condition for determinacy.
\end{nuance}

\subsection*{Cumulants}

The \emph{cumulant generating function} is $K_X(t)=\ln M_X(t)$, and its
derivatives at zero are the \emph{cumulants} $\kappa_1,\kappa_2,\dots$ Because
$M_{X+Y}=M_XM_Y$ for independent variables, $K_{X+Y}=K_X+K_Y$, so cumulants are
additive. The first two are the mean and variance, $\kappa_3$ is the third central
moment. For the normal, $K(t)=\mu t+\sigma^2t^2/2$, so all cumulants beyond the
second vanish; for the Poisson, $K(t)=\lambda(e^t-1)$, so every cumulant equals
$\lambda$.

\section{Probability generating functions}

For a random variable $N$ taking values in $\{0,1,2,\dots\}$ with
$p_k=\Prob(N=k)$, the \emph{probability generating function} (PGF) is
\[
G_N(s)=\Exp[s^N]=\sum_{k=0}^\infty p_ks^k,\qquad\abs s\le1.
\]
The series converges absolutely for $\abs s\le1$ since $\sum p_k=1$.

\begin{theorem}[Properties of the PGF]\label{thm:pgf}
\begin{enumerate}[label=(\roman*)]
\item $G_N(1)=1$ and $G_N$ determines the PMF: $p_k=G_N^{(k)}(0)/k!$.
\item $\Exp[N]=G_N'(1^-)$ and $\Exp[N(N-1)]=G_N''(1^-)$ (derivatives from the left;
both sides may be $+\infty$), so $\Var(N)=G''(1)+G'(1)-G'(1)^2$.
\item If $N,M$ are independent, $G_{N+M}=G_NG_M$.
\end{enumerate}
\end{theorem}
\begin{proof}
(i) A power series is determined by its coefficients, and the Taylor coefficients
are $G^{(k)}(0)/k!$. (ii) Differentiating the power series term by term for
$\abs s<1$ gives $G'(s)=\sum kp_ks^{k-1}$ and $G''(s)=\sum k(k-1)p_ks^{k-2}$; let
$s\uparrow1$ using monotone convergence (all terms are non-negative). (iii)
$s^{N+M}=s^Ns^M$ with independent factors.
\end{proof}

\begin{theorem}[Random sums: Wald's identity and compound variance]\label{thm:randomsum}
Let $X_1,X_2,\dots$ be i.i.d.\ non-negative integer-valued with PGF $G_X$, and let $N$ be
independent of the $X_i$, also non-negative integer-valued with PGF $G_N$. Then
$S=\sum_{i=1}^NX_i$ has PGF $G_S(s)=G_N\bigl(G_X(s)\bigr)$. If the means
and variances are finite,
\[
\Exp[S]=\Exp[N]\Exp[X],\qquad\Var(S)=\Exp[N]\Var(X)+\Var(N)\,(\Exp[X])^2.
\]
\end{theorem}
\begin{proof}
Condition on $N$ (law of total probability for expectations over the partition
$\{N=n\}$):
\[
G_S(s)=\sum_n\Prob(N=n)\,\Exp\bigl[s^{X_1+\dots+X_n}\bigr]=\sum_n\Prob(N=n)G_X(s)^n=G_N(G_X(s)).
\]
Differentiating, $G_S'(s)=G_N'(G_X(s))G_X'(s)$, so $G_S'(1)=G_N'(1)G_X'(1)$ (as
$G_X(1)=1$), which is the mean formula. Differentiating again,
$G_S''(1)=G_N''(1)G_X'(1)^2+G_N'(1)G_X''(1)$, and substituting into
$\Var(S)=G_S''(1)+G_S'(1)-G_S'(1)^2$ and simplifying with $G''(1)=\Exp[X(X-1)]$
gives the variance formula.
\end{proof}

\begin{example}[Compound Poisson thinning]
Each of $N\sim\mathrm{Poisson}(\lambda)$ customers independently buys with
probability $p$. The number of buyers is $S=\sum_{i\le N}X_i$ with
$X_i\sim\mathrm{Bernoulli}(p)$, $G_X(s)=q+ps$. Then
$G_S(s)=\exp\bigl(\lambda(q+ps-1)\bigr)=\exp\bigl(\lambda p(s-1)\bigr)$, which is the
Poisson$(\lambda p)$ PGF. So \emph{thinning a Poisson count gives a Poisson count}.
\end{example}

\subsection*{Branching processes and extinction}

\begin{example}[Galton--Watson extinction]\label{ex:branching}
Start with one individual. Each individual independently has $\xi$ children, with
PMF $p_k$ and PGF $G$. Let $Z_n$ be the size of generation $n$ and
$\pi=\Prob(\text{eventual extinction})$. By conditioning on the first generation: if
the root has $k$ children, each of the $k$ independent subtrees dies out with
probability $\pi$, so
\[
\pi=\sum_kp_k\pi^k=G(\pi).
\]
\begin{theorem}\label{thm:extinction}
Assume $p_1<1$. The extinction probability is the \emph{smallest} root of
$s=G(s)$ in $[0,1]$. If $\Exp[\xi]=G'(1)\le1$ then $\pi=1$; if $\Exp[\xi]>1$ then $\pi<1$.
\end{theorem}
\begin{proof}
Let $\pi_n=\Prob(Z_n=0)$. By the same conditioning, $\pi_{n+1}=G(\pi_n)$ with
$\pi_0=0$. The events $\{Z_n=0\}$ increase to the extinction event, so
$\pi_n\uparrow\pi$ by continuity (Theorem~\ref{thm:continuity}), and since $G$ is
continuous, $\pi=G(\pi)$. If $s^*$ is any root of $s=G(s)$ in $[0,1]$, then
$\pi_0=0\le s^*$ and, $G$ being non-decreasing, $\pi_{n+1}=G(\pi_n)\le G(s^*)=s^*$ by
induction; hence $\pi\le s^*$, so $\pi$ is the smallest root. $G$ is convex on $[0,1]$
(non-negative coefficients), $G(1)=1$. If $G'(1)\le1$ the graph of $G$ lies above the
diagonal on $[0,1)$ (convexity plus slope at most $1$ at $s=1$, using $p_1<1$ for
strictness), so the only root is $1$. If $G'(1)>1$, then $G(s)<s$ for $s$ slightly
below $1$, while $G(0)\ge0$, so by continuity there is a root in $[0,1)$.
\end{proof}
For instance with offspring law $p_0=\tfrac14,p_1=\tfrac14,p_2=\tfrac12$: the mean is
$\tfrac54>1$ and $s=\tfrac14+\tfrac14s+\tfrac12s^2$ gives
$2s^2-3s+1=0$, i.e.\ $s=\tfrac12$ or $s=1$. So $\pi=\tfrac12$.
\end{example}

\section{Characteristic functions}\label{sec:cf}

\begin{definition}[Characteristic function]
The \emph{characteristic function} of $X$ is
$\varphi_X(t)=\Exp[e^{itX}]=\Exp[\cos tX]+i\,\Exp[\sin tX]$, for $t\in\R$.
\end{definition}

Unlike the MGF it \emph{always exists}, because $\abs{e^{itX}}=1$.

\begin{theorem}[Properties]\label{thm:cfprops}
\begin{enumerate}[label=(\roman*)]
\item $\varphi_X(0)=1$, $\abs{\varphi_X(t)}\le1$, and $\varphi_X(-t)=\overline{\varphi_X(t)}$.
\item $\varphi_{aX+b}(t)=e^{itb}\varphi_X(at)$.
\item If $X,Y$ are independent, $\varphi_{X+Y}=\varphi_X\varphi_Y$.
\item (Moments, cited) If $\Exp\abs X^k<\infty$ then $\varphi_X$ is $k$ times
continuously differentiable and $\varphi_X^{(k)}(0)=i^k\Exp[X^k]$. In particular,
$\varphi_X(t)=1+it\Exp[X]-\tfrac12t^2\Exp[X^2]+o(t^2)$ when $\Exp[X^2]<\infty$.
\item (Uniqueness, cited) $\varphi_X=\varphi_Y$ implies $X\overset d=Y$.
\item (L\'evy continuity theorem, cited; proved in Chapter~\ref{ch:measure})
$X_n\xrightarrow{d}X$ iff $\varphi_{X_n}(t)\to\varphi_X(t)$ for every $t$ (where the
limit function is continuous at $0$).
\end{enumerate}
\end{theorem}
\begin{proof}[Proof of (i)--(iii)]
(i) $\abs{\Exp[e^{itX}]}\le\Exp\abs{e^{itX}}=1$; conjugation: $\overline{e^{itX}}=e^{-itX}$.
(ii) $e^{it(aX+b)}=e^{itb}e^{i(at)X}$. (iii) As for MGFs, $e^{itX}$ and $e^{itY}$ are
independent and bounded.
\end{proof}

\begin{example}[Standard normal]\label{ex:cfnormal}
$\varphi_Z(t)=e^{-t^2/2}$. \emph{Derivation by differential equation.} Differentiating
under the integral (justified since $\abs{xe^{itx}\varphi(x)}$ is integrable),
\[
\varphi_Z'(t)=\int ix e^{itx}\varphi(x)\,dx=-i\int e^{itx}\varphi'(x)\,dx
=-i\Bigl([e^{itx}\varphi(x)]_{-\infty}^\infty-it\int e^{itx}\varphi(x)\,dx\Bigr)=-t\varphi_Z(t),
\]
using $x\varphi=-\varphi'$ and integration by parts. With $\varphi_Z(0)=1$ the unique
solution is $e^{-t^2/2}$. Hence $\varphi_{N(\mu,\sigma^2)}(t)=e^{i\mu t-\sigma^2t^2/2}$.
\end{example}

\begin{example}[Cauchy: the law of large numbers fails]\label{ex:cauchycf}
The standard Cauchy has $\varphi(t)=e^{-\abs t}$. \emph{Derivation.} Start from the
two-sided exponential density $\tfrac12e^{-\abs x}$, whose characteristic function is
$\tfrac12\int e^{-\abs x}e^{itx}dx=\tfrac12\bigl(\tfrac1{1-it}+\tfrac1{1+it}\bigr)=\tfrac1{1+t^2}$.
The Fourier inversion theorem (cited here, proved in Chapter~\ref{ch:measure}) says that
when $\varphi$ is integrable the density is $f(x)=\frac1{2\pi}\int\varphi(t)e^{-itx}dt$.
Applying this with $\varphi(t)=\frac1{1+t^2}$ and $f(x)=\tfrac12e^{-\abs x}$ gives
$\tfrac12e^{-\abs x}=\frac1{2\pi}\int\frac{e^{-itx}}{1+t^2}dt$, i.e.\
$\int\frac{e^{itx}}{\pi(1+t^2)}\,dt=e^{-\abs x}$. Read with the roles of $t$ and $x$
swapped, this says exactly that the Cauchy density $\frac1{\pi(1+x^2)}$ has
characteristic function $e^{-\abs t}$. If $X_1,\dots,X_n$ are i.i.d.\ standard
Cauchy, the sample mean has
\[
\varphi_{\bar X_n}(t)=\Bigl(\varphi\bigl(\tfrac tn\bigr)\Bigr)^n=\bigl(e^{-\abs t/n}\bigr)^n=e^{-\abs t},
\]
so $\bar X_n$ is again standard Cauchy for every $n$. Averaging does \emph{nothing}:
the distribution of $\bar X_n$ never concentrates. The law of large numbers (next
chapters) requires a finite mean. This also gives a short proof of the stability
claim in Exercise 5 of Chapter~\ref{ch:jointcont}.
\end{example}

\begin{interview}{which transform to use}
Use the \textbf{PGF} for non-negative integer variables, especially for random
sums and branching processes; the \textbf{MGF} for tail bounds (Chernoff,
Chapter~\ref{ch:condexp}) and quick moment computations when it exists; the
\textbf{characteristic function} when you need existence for every distribution
(limit theorems). All three satisfy ``transform of a sum of independents equals
product of transforms''.
\end{interview}

\section*{Exercises}
\addcontentsline{toc}{section}{Exercises}

\begin{enumerate}
\item Compute the MGF of the geometric distribution and use it to rederive
$\Exp[X]=1/p$ and $\Var(X)=q/p^2$.
\item Show that if $X\sim\mathrm{Bin}(n,p)$ then $G_X(s)=(q+ps)^n$, and use
Theorem~\ref{thm:pgf} to recover the mean and variance.
\item Find the distribution of $N_1+N_2$ if $N_i$ are independent Geom$(p)$ using
PGFs.
\item Let $X\sim N(0,1)$. Show that $\Exp[e^{tX^2}]=(1-2t)^{-1/2}$ for $t<\frac12$ and
deduce that $X^2\sim\mathrm{Gamma}(\tfrac12,\tfrac12)$.
\item A fair die is rolled; if the result is $N$, then $N$ fair coins are tossed.
Find the mean and variance of the number of heads using
Theorem~\ref{thm:randomsum}.
\item Compute the extinction probability of a branching process with offspring
distribution Poisson$(2)$ numerically (solve $s=e^{2(s-1)}$).
\item Prove: if $X$ has an MGF finite near $0$ and $\Exp[X]=0$, then
$M_X(t)=1+\tfrac12\Var(X)t^2+O(t^3)$.
\item Show that $\varphi_X$ is real-valued if and only if $X$ and $-X$ have the same
distribution.
\item (Challenge) Prove that the characteristic function is uniformly continuous
on $\R$ (Hint: $\abs{\varphi(t+h)-\varphi(t)}\le\Exp\abs{e^{ihX}-1}$ and
dominated convergence).
\end{enumerate}
